\documentclass[UTF8,12pt]{ctexart}
\usepackage[a4paper,margin=24mm]{geometry}
\usepackage{amsmath,amssymb,bm,graphicx,adjustbox}
\newcommand{\fitmath}[1]{\begin{adjustbox}{max width=\linewidth}$\displaystyle #1$\end{adjustbox}}
\setlength{\parindent}{0pt}
\setlength{\parskip}{.7em}
\sloppy
\allowdisplaybreaks
\begin{document}
\section*{1995 年考研数学三 · 真题}
\subsection*{1995 年全国硕士研究生招生考试试题（试卷 IV）}

\subsection*{一、填空题（本题共 5 小题，每小题 3 分，满分 15 分）}

（1）设 \(\displaystyle f(x)=\dfrac{\displaystyle 1-x}{\displaystyle 1+x}\)，则 \(\displaystyle f^{(n)}(x)=\)\_\_\_\_\_\_。


（2）设 \(\displaystyle z=xyf\left(\dfrac{\displaystyle y}{\displaystyle x}\right)\)，\(\displaystyle f(u)\) 可导，则 \(\displaystyle xz_x^{\prime}+yz_y^{\prime}=\)\_\_\_\_\_\_。


（3）设 \(\displaystyle f^{\prime}(\ln x)=1+x\)，则 \(\displaystyle f(x)=\)\_\_\_\_\_\_。


（4）设 \(\displaystyle A=\begin{pmatrix}1&0&0\\2&2&0\\3&4&5\end{pmatrix}\)，\(\displaystyle A^*\) 是 \(\displaystyle A\) 的伴随矩阵，则 \(\displaystyle (A^*)^{-1}=\)\_\_\_\_\_\_。


（5）（超纲题）设 \(\displaystyle X_1,\allowbreak \ldots,\allowbreak X_n\) 是来自正态总体 \(\displaystyle N(\mu,\allowbreak \sigma^2)\) 的简单随机样本，其中参数 \(\displaystyle \mu\) 和 \(\displaystyle \sigma^2\) 未知，记 \(\displaystyle \bar X=\dfrac{\displaystyle 1}{\displaystyle n}\displaystyle \sum_{i=1}^nX_i,\allowbreak \ Q^2=\displaystyle \sum_{i=1}^n(X_i-\bar X)^2\)，则假设 \(\displaystyle H_0:\mu=0\) 的 \(\displaystyle t\) 检验使用统计量 \(\displaystyle t=\)\_\_\_\_\_\_。


\subsection*{二、选择题（本题共 5 小题，每小题 3 分，满分 15 分）}

（1）设 \(\displaystyle f(x)\) 为可导函数，且满足条件 \(\displaystyle \lim_{x\to0}\dfrac{\displaystyle f(1)-f(1-x)}{\displaystyle 2x}=-1\)，则曲线 \(\displaystyle y=f(x)\) 在点 \(\displaystyle (1,\allowbreak f(1))\) 处的切线斜率为（　）。


（A）2。 （B）\(\displaystyle -1\)。 （C）\(\displaystyle \dfrac{\displaystyle 1}{\displaystyle 2}\)。 （D）\(\displaystyle -2\)。


（2）下列广义积分发散的是（　）。


（A）\(\displaystyle \int_{-1}^1\dfrac{\displaystyle 1}{\displaystyle \sin x}\,dx\)。 （B）\(\displaystyle \int_{-1}^1\dfrac{\displaystyle 1}{\displaystyle \sqrt{1-x^2}}\,dx\)。 （C）\(\displaystyle \int_0^{+\infty}e^{-x^2}\,dx\)。 （D）\(\displaystyle \int_2^{+\infty}\dfrac{\displaystyle 1}{\displaystyle x\ln^2x}\,dx\)。


（3）设矩阵 \(\displaystyle A_{m\times n}\) 的秩为 \(\displaystyle r(A)=m<n\)，\(\displaystyle E_m\) 为 \(\displaystyle m\) 阶单位矩阵，则下述结论中正确的是（　）。


（A）\(\displaystyle A\) 的任意 \(\displaystyle m\) 个列向量必线性无关。 （B）\(\displaystyle A\) 的任意一个 \(\displaystyle m\) 阶子式不等于零。 （C）若矩阵 \(\displaystyle B\) 满足 \(\displaystyle BA=O\)，则 \(\displaystyle B=O\)。 （D）\(\displaystyle A\) 通过初等行变换，必可以化为 \(\displaystyle (E_m,\allowbreak O)\) 的形式。


（4）设随机变量 \(\displaystyle X\) 和 \(\displaystyle Y\) 独立同分布，记 \(\displaystyle U=X-Y,\allowbreak \ V=X+Y\)，则随机变量 \(\displaystyle U\) 与 \(\displaystyle V\) 必然（　）。


（A）不独立。 （B）独立。 （C）相关系数不为零。 （D）相关系数为零。


（5）设随机变量 \(\displaystyle X\sim N(\mu,\allowbreak \sigma^2)\)，则随着 \(\displaystyle \sigma\) 的增大，概率 \(\displaystyle P\{|X-\mu|<\sigma\}\)（　）。


（A）单调增大。 （B）单调减小。 （C）保持不变。 （D）增减不定。


\subsection*{试卷 IV（续）}

三、（本题满分 6 分）设 \(\displaystyle f(x)=\begin{cases}\dfrac{\displaystyle 2}{\displaystyle x^2}(1-\cos x),\allowbreak &x<0,\allowbreak \\1,\allowbreak &x=0,\allowbreak \\\dfrac{\displaystyle 1}{\displaystyle x}\displaystyle \int_0^x\cos(t^2)\,dt,\allowbreak &x>0,\allowbreak \end{cases}\) 试讨论 \(\displaystyle f(x)\) 在 \(\displaystyle x=0\) 处的连续性和可导性。


四、（本题满分 6 分）已知连续函数 \(\displaystyle f(x)\) 满足条件 \(\displaystyle f(x)=\displaystyle \int_0^{3x}f\left(\dfrac{\displaystyle t}{\displaystyle 3}\right)\,dt+e^{2x}\)，求 \(\displaystyle f(x)\)。


五、（本题满分 6 分）将函数 \(\displaystyle y=\ln(1-x-2x^2)\) 展成 \(\displaystyle x\) 的幂级数，并指出其收敛区间。


六、（本题满分 5 分）计算二次积分 \(\displaystyle I=\displaystyle \int_{-\infty}^{+\infty}\displaystyle \int_{-\infty}^{+\infty}\displaystyle \min\{x,\allowbreak y\}e^{-(x^2+y^2)}\,dx\,dy\)。


七、（本题满分 6 分）设某产品的需求函数为 \(\displaystyle Q=Q(P)\)，收益函数为 \(\displaystyle R=PQ\)，其中 \(\displaystyle P\) 为产品价格，\(\displaystyle Q\) 为需求量（产品的产量），\(\displaystyle Q(P)\) 是单调减函数。如果当价格为 \(\displaystyle P_0\)，对应产量为 \(\displaystyle Q_0\) 时，边际收益 \(\displaystyle \left.\dfrac{\displaystyle dR}{\displaystyle dQ}\right|_{Q=Q_0}=a>0\)，收益对价格的边际效应 \(\displaystyle \left.\dfrac{\displaystyle dR}{\displaystyle dP}\right|_{P=P_0}=c<0\)，需求对价格的弹性 \(\displaystyle E_P=b>1\)，求 \(\displaystyle P_0\) 和 \(\displaystyle Q_0\)。


八、（本题满分 6 分）设 \(\displaystyle f(x),\allowbreak g(x)\) 在区间 \(\displaystyle [-a,\allowbreak a]\)（\(\displaystyle a>0\)）上连续，\(\displaystyle g(x)\) 为偶函数，且 \(\displaystyle f(x)\) 满足条件 \(\displaystyle f(x)+f(-x)=A\)（\(\displaystyle A\) 为常数）。（1）证明 \(\displaystyle \int_{-a}^a f(x)g(x)\,dx=A\displaystyle \int_0^a g(x)\,dx\)；（2）利用（1）的结论计算定积分 \(\displaystyle \int_{-\pi/2}^{\pi/2}|\sin x|\arctan e^x\,dx\)。


九、（本题满分 9 分）已知向量组（I）\(\displaystyle \alpha_1,\allowbreak \alpha_2,\allowbreak \alpha_3\)；（II）\(\displaystyle \alpha_1,\allowbreak \alpha_2,\allowbreak \alpha_3,\allowbreak \alpha_4\)；（III）\(\displaystyle \alpha_1,\allowbreak \alpha_2,\allowbreak \alpha_3,\allowbreak \alpha_5\)。如果各向量组的秩分别为 \(\displaystyle r(\mathrm I)=r(\mathrm{II})=3,\allowbreak  r(\mathrm{III})=4\)。证明：向量组 \(\displaystyle \alpha_1,\allowbreak \alpha_2,\allowbreak \alpha_3,\allowbreak \alpha_5-\alpha_4\) 的秩为 4。


十、（本题满分 10 分）已知二次型 \(\displaystyle f(x_1,\allowbreak x_2,\allowbreak x_3)=4x_2^2-3x_3^2+4x_1x_2-4x_1x_3+8x_2x_3\)。（1）写出二次型 \(\displaystyle f\) 的矩阵表达式；


\subsection*{试卷 IV（续）}

十、（2）用正交变换把二次型 f 化为标准形，并写出相应的正交矩阵。


十一、（本题满分 8 分）假设一家生产的每台仪器，以概率 0.70 可以直接出厂；以概率 0.30 需进一步调试，经调试后以概率 0.80 可以出厂；以概率 0.20 定为不合格品不能出厂。现该厂生产了 \(\displaystyle n\)（\(\displaystyle n\geq2\)）台仪器（假设各台仪器的生产过程相互独立）。求：


（1）全部能出厂的概率 \(\displaystyle \alpha\)；（2）其中恰好有两台不能出厂的概率 \(\displaystyle \beta\)；（3）其中至少有两台不能出厂的概率 \(\displaystyle \theta\)。


十二、（本题满分 8 分）已知随机变量 \(\displaystyle X\) 和 \(\displaystyle Y\) 的联合概率密度为 \(\displaystyle \varphi(x,\allowbreak y)=\begin{cases}4xy,\allowbreak &0\leq x\leq1,\allowbreak \ 0\leq y\leq1,\allowbreak \\0,\allowbreak &\text{其他},\allowbreak \end{cases}\) 求 \(\displaystyle X\) 和 \(\displaystyle Y\) 的联合分布函数 \(\displaystyle F(x,\allowbreak y)\)。


\subsection*{1995 年全国硕士研究生招生考试试题（试卷 V）}

\subsection*{一、填空题（本题共 5 小题，每小题 3 分，满分 15 分）}

（1）\(\displaystyle \lim_{x\to\infty}\left(\dfrac{\displaystyle 1+x}{\displaystyle x}\right)^{ax}=\displaystyle \int_{-\infty}^a te^t\,dt\)，则常数 \(\displaystyle a=\)\_\_\_\_\_\_。


（2）【同试卷 IV 第一、（2）题】


（3）【同试卷 IV 第一、（3）题】


（4）【同试卷 IV 第一、（4）题】


（5）设 \(\displaystyle X\) 是一个随机变量，其概率密度为 \(\displaystyle f(x)=\begin{cases}1+x,\allowbreak &-1\leq x\leq0,\allowbreak \\1-x,\allowbreak &0<x\leq1,\allowbreak \\0,\allowbreak &\text{其他},\allowbreak \end{cases}\) 则方差 \(\displaystyle D(X)=\)\_\_\_\_\_\_。


\subsection*{二、选择题（本题共 5 小题，每小题 3 分，满分 15 分）}

（1）【同试卷 IV 第二、（1）题】


（2）【同试卷 IV 第二、（2）题】


（3）设 \(\displaystyle n\) 维行向量 \(\displaystyle \alpha=(\dfrac{\displaystyle 1}{\displaystyle 2},\allowbreak 0,\allowbreak \ldots,\allowbreak 0,\allowbreak \dfrac{\displaystyle 1}{\displaystyle 2})\)，矩阵 \(\displaystyle A=E-\alpha^{\mathrm T}\alpha,\allowbreak \ B=E+2\alpha^{\mathrm T}\alpha\)，其中 \(\displaystyle E\) 为 \(\displaystyle n\) 阶单位矩阵，则 \(\displaystyle AB\) 等于（　）。


（A）\(\displaystyle O\)。 （B）\(\displaystyle -E\)。 （C）\(\displaystyle E\)。 （D）\(\displaystyle E+\alpha^{\mathrm T}\alpha\)。


（4）设矩阵 \(\displaystyle A_{m\times n}\) 的秩为 \(\displaystyle r(A)=m<n\)，\(\displaystyle E_m\) 为 \(\displaystyle m\) 阶单位矩阵，下述结论中正确的是（　）。


（A）\(\displaystyle A\) 的任意 \(\displaystyle m\) 个列向量必线性无关。 （B）\(\displaystyle A\) 的任意一个 \(\displaystyle m\) 阶子式不等于零。 （C）非齐次线性方程组 \(\displaystyle Ax=b\) 一定有无穷多解。 （D）\(\displaystyle A\) 通过初等行变换，必可以化为 \(\displaystyle (E_m,\allowbreak O)\) 的形式。


\subsection*{试卷 V（续）}

二、（5）【同试卷 IV 第二、（5）题】


三、（本题满分 6 分）【同试卷 IV 第三题】


四、（本题满分 6 分）求不定积分 \(\displaystyle \int(\arcsin x)^2\,dx\)。


五、（本题满分 7 分）【同试卷 IV 第八题】


六、（本题满分 6 分）【同试卷 IV 第七题】


七、（本题满分 5 分）设 \(\displaystyle f(x)\) 在区间 \(\displaystyle [a,\allowbreak b]\) 上连续，在 \(\displaystyle (a,\allowbreak b)\) 内可导，证明：在 \(\displaystyle (a,\allowbreak b)\) 内至少存在一点 \(\displaystyle \xi\)，使 \(\displaystyle \dfrac{\displaystyle bf(b)-af(a)}{\displaystyle b-a}=f(\xi)+\xi f^{\prime}(\xi)\)。


八、（本题满分 9 分）求二元函数 \(\displaystyle z=f(x,\allowbreak y)=x^2y(4-x-y)\) 在由直线 \(\displaystyle x+y=6\)，\(\displaystyle x\) 轴和 \(\displaystyle y\) 轴所围成的闭区域 \(\displaystyle D\) 上的极值、最大值与最小值。


九、（本题满分 8 分）对于线性方程组 \(\displaystyle \begin{cases}\lambda x_1+x_2+x_3=\lambda-3,\allowbreak \\x_1+\lambda x_2+x_3=-2,\allowbreak \\x_1+x_2+\lambda x_3=-2,\allowbreak \end{cases}\) 讨论 \(\displaystyle \lambda\) 取何值时，方程组无解、有唯一解和无穷多组解。在方程组有无穷多组解时，试用其导出组的基础解系表示全部解。


十、（本题满分 8 分）设 3 阶矩阵 \(\displaystyle A\) 满足 \(\displaystyle A\alpha_i=i\alpha_i\)（\(\displaystyle i=1,\allowbreak 2,\allowbreak 3\)），其中列向量 \(\displaystyle \alpha_1=(1,\allowbreak 2,\allowbreak 2)^{\mathrm T},\allowbreak \ \alpha_2=(2,\allowbreak -2,\allowbreak 1)^{\mathrm T},\allowbreak \ \alpha_3=(-2,\allowbreak -1,\allowbreak 2)^{\mathrm T}\)。试求矩阵 \(\displaystyle A\)。


十一、（本题满分 8 分）【同试卷 IV 第十一题】


十二、（本题满分 7 分）假设随机变量 \(\displaystyle X\) 服从参数为 2 的指数分布，证明：\(\displaystyle Y=1-e^{-2X}\) 在区间 \(\displaystyle (0,\allowbreak 1)\) 上服从均匀分布。

\end{document}
